\subsubsection*{EXP 4: Algebraic subgroup analysis}

We checked whether the learned matrices $P$ correspond to known geometric symmetries (flip, rot90, rot180, rot270) or their subgroups.

\textbf{Distances to known symmetries:}
\begin{table}[h]
\centering
\caption{EXP~4: Frobenius distances of learned $P$ to geometric symmetries.}
\begin{tabular}{lccc}
\hline
 & $P_1$ & $P_2$ & $P_3$ \\
\hline
Identity & 10.89 & 10.87 & 10.75 \\
H-Flip   & 10.81 & 10.80 & 10.61 \\
\textbf{V-Flip} & \textbf{10.53} & \textbf{10.49} & \textbf{10.39} \\
Rot90    & 10.69 & 10.68 & 10.55 \\
Rot180   & 10.73 & 10.70 & 10.59 \\
Rot270   & 10.89 & 10.87 & 10.75 \\
\hline
\end{tabular}
\end{table}

No power $P^2$ through $P^6$ approached the identity (distances 8.9--10.3). No matrix exhibited the conjugation relation $PJP^\top \approx J$ with any known symmetry.

\textbf{Conclusion:} The model \textbf{does not recover} standard geometric symmetries. The closest is V-Flip (10.39), which is consistent with Task~B containing vflip, but the distance is still enormous. The model learns \textbf{functional} permutations (minimizing loss), not \textbf{geometric} ones (corresponding to physical transformations).
